\documentclass{beamer}
\usepackage[spanish]{babel}
\usepackage{amsmath, amssymb, graphicx, xcolor, multirow, diagbox, tikz}
\usetikzlibrary{positioning}
\usetheme{default}

\title{Clase 1: Introducción a la Toma de Decisiones Secuenciales}
\subtitle{Métodos para la Gestión - Maestría en Negocios}
\author{Daniela Pucci de Farias}
\institute{Universidad Torcuato Di Tella}
\date{Junio 2025}

\begin{document}

\frame{\titlepage}

\begin{frame}{El Corazón de la Gestión: Decisiones en el Tiempo}
\begin{itemize}
    \item La mayoría de los problemas de negocios no son eventos únicos.
    \item \textbf{Ejemplos Reales:}
    \begin{itemize}
        \item Precios dinámicos (ajustar según demanda).
        \item Gestión de inventario (pedir hoy para vender mañana).
        \item Inversión de capital (gastar ahora para crecer después).
    \end{itemize}
    \item \textbf{Incertidumbre:} El futuro no es determinístico; es probabilístico.
\end{itemize}
\end{frame}

\begin{frame}{Intuición de RL: Exploración vs. Explotación}
\begin{itemize}
    \item Antes de la matemática, pensemos en el dilema del manager:
    \item \textbf{Explotación:} Usar la estrategia que funcionó el mes pasado para asegurar ventas conocidas.
    \item \textbf{Exploración:} Probar un nuevo canal de marketing o un nuevo precio para ver si hay una oportunidad mejor.
    \item \textbf{Gestión:} ¿Cómo equilibramos "ir a lo seguro" con "aprender del mercado"?
\end{itemize}
\end{frame}

\begin{frame}{El Marco: Agente y Ambiente}
\centering
\begin{tikzpicture}[node distance=2cm, auto]
    \node [rectangle, draw, minimum width=2cm] (agent) {Agente (Manager)};
    \node [rectangle, draw, minimum width=2cm, below of=agent] (env) {Ambiente (Mercado)};
    \draw [->] (agent.east) -- ++(1,0) |- node [near start] {Acción $a_t$} (env.east);
    \draw [->] (env.west) -- ++(-1,0) |- node [near start] {Estado $s_{t+1}$, Recompensa $r_t$} (agent.west);
\end{tikzpicture}
\begin{itemize}
    \item \textbf{Estado ($s$):} Resumen de la situación actual.
    \item \textbf{Acción ($a$):} Lo que decidimos hacer.
    \item \textbf{Recompensa ($r$):} Feedback (ganancia, costo, satisfacción).
\end{itemize}
\end{frame}

\begin{frame}{Caso de Estudio: Participación de Mercado en Ghana}
\begin{itemize}
    \item Modelamos la lealtad del cliente como una Cadena de Markov.
    \item \textbf{Datos del Problema:}
    \begin{itemize}
        \item Operador M: 71.6\% de retención.
        \item Operador V: 38.7\% de retención.
    \end{itemize}
    \item \textbf{Pregunta de Negocio:} Si hoy tengo el 33\% del mercado, ¿dónde estaré en 2 años?
    \item \textbf{Simulación:} Veremos en Python cómo la "Distribución Estacionaria" nos dice el destino final del mercado independientemente de hoy.
\end{itemize}
\end{frame}

\begin{frame}{La Propiedad de Markov: "Lo que importa es hoy"}
\begin{block}{Definición}
$P(s_{t+1} | s_t, s_{t-1}, ..., s_0) = P(s_{t+1} | s_t)$
\end{block}
\begin{itemize}
    \item Para un manager, esto significa elegir un \textbf{Estado} que sea un "resumen suficiente".
    \item Si el inventario de ayer afecta la venta de mañana de una forma que el inventario de hoy no explica, ¡tu estado está incompleto!
\end{itemize}
\end{frame}

\begin{frame}{Laboratorio 1: Introducción a Python para MDPs}
\begin{itemize}
    \item Entorno: Google Colab.
    \item \textbf{Ejercicio 1:} Definir matrices de transición para el mercado de telefonía.
    \item \textbf{Ejercicio 2:} Proyectar la cuota de mercado a 12 meses.
    \item \textbf{Visualización:} Graficar la convergencia hacia el equilibrio.
\end{itemize}
\end{frame}

\end{document}
